% ID: SEL-D0044-OB-E04
% ESTADO: validada. S045, cotejada con examen y criterios.
\textbf{a)} Sin reemplazo y sin orden, las muestras son
$(6,9)$, $(6,12)$ y $(9,12)$, cada una con probabilidad $1/3$.
Sus medias son $7.5$, $9$ y $10.5$, con esperanza 9.
Por tanto la varianza de la distribución de las medias es
\[
\frac{(7.5-9)^2+(9-9)^2+(10.5-9)^2}{3}=\boxed{1.5}.
\]
Si se distingue el orden sin reemplazo, cada media aparece
dos veces y la varianza es la misma. Con reemplazo hay
nueve pares ordenados equiprobables. La varianza poblacional
es $(9+0+9)/3=6$ y la de la media de dos extracciones
independientes es $6/2=3$.

\textbf{b)} La producción total es $40+15+25+120=200$.
La fracción común es $40/200=1/5$. Se toma un quinto
de cada producción: $\boxed{8,\ 3,\ 5,\ 24}$ unidades
de A, B, C y D, respectivamente. Suman 40.
